
\documentclass{article}
\usepackage{colortbl}
\usepackage{makecell}
\usepackage{multirow}
\usepackage{supertabular}

\begin{document}

\newcounter{utterance}

\centering \large Interaction Transcript for game `codenames', experiment `association\_difficult', episode 1 with glm{-}5.2.
\vspace{24pt}

{ \footnotesize  \setcounter{utterance}{1}
\setlength{\tabcolsep}{0pt}
\begin{supertabular}{c@{$\;$}|p{.15\linewidth}@{}p{.15\linewidth}p{.15\linewidth}p{.15\linewidth}p{.15\linewidth}p{.15\linewidth}}
    \# & $\;$A & \multicolumn{4}{c}{Game Master} & $\;\:$B\\
    \hline

    \theutterance \stepcounter{utterance}  
    & & & \multicolumn{2}{p{0.3\linewidth}}{
        \cellcolor[rgb]{0.9,0.9,0.9}{
            \makecell[{{p{\linewidth}}}]{
                \texttt{\tiny{[GM$|$GM]}}
                {'hidden': {'team': ['gümüş', 'iris', 'etek', 'yurt', 'şeker', 'adam kaçırma', 'tipi', 't-shirt', 'nergis'], 'innocent': ['tavan', 'tüberküloz', 'leylak', 'iglo', 'oda', 'kot pantolon', 'veba'], 'opponent': ['duvar', 'petunya', 'beyaz', 'darp', 'mağara', 'hardal', 'çatı', 'külot'], 'assassin': ['hız yapma']}, 'revealed': {'team': {'team': [], 'innocent': [], 'opponent': [], 'assassin': []}, 'opponent': {'team': [], 'innocent': [], 'opponent': [], 'assassin': []}}}
            }
        }
    }
    & & \\ \\

    \theutterance \stepcounter{utterance}  
    & & \multicolumn{4}{p{0.6\linewidth}}{
        \cellcolor[rgb]{0.9,0.9,0.9}{
            \makecell[{{p{\linewidth}}}]{
                \texttt{\tiny{[P1$\langle$GM]}}
                \texttt{Aşağıdaki listeden takımının kelimelerinden bazılarıyla ilişkili tek bir tek kelimelik ipucu ver. İpucunun kaç kelimeyle ve hangi kelimelerle ilişkili olacağını sen seçebilirsin.} \\
\texttt{Amaç, diğer takım kendi kelimelerinin hepsini bulmadan önce, verdiğin ipuçlarına dayanarak takım arkadaşının senin takımına ait tüm kelimeleri tahmin etmesini sağlamak. Takım arkadaşın bir "rakip" ya da bir "dikkat dağıtıcı" kelime tahmin ederse, tahmin sırası biter ve sıra rakip takıma geçer.  Eğer partnerin bir "öldüren" kelime tahmin ederse, oyun biter ve takımın anında kaybeder.} \\
\texttt{İpucun ne kadar çok kelimeyle ilişkili olursa oyunu o kadar hızlı kazanabilirsin, ancak çağrışımların takım arkadaşın için tahmin etmesi o kadar da zor olabilir; bu yüzden ipucu kelimeni akıllıca seç. İpucu kelimesi hedef kelimelerle anlamsal olarak ilişkili olmalıdır; listelerdeki kelimelerden biri olamaz ve onların parçalarını içeremez.} \\
\\ 
\texttt{Tek kelimelik ipucunu ve ilgili hedef kelimelerin virgülle ayrılmış listesini her zaman aşağıdaki biçimde ver ve yanıtlarını olabildiğince kısa tut; bu formda gereken dışında asla başka metin ekleme:} \\
\\ 
\texttt{İPUCU: <KELİME>} \\
\texttt{HEDEFLER: <HEDEFLER>} \\
\\ 
\texttt{Takımının kelimeleri: gümüş, iris, etek, yurt, şeker, adam kaçırma, tipi, t{-}shirt, nergis.} \\
\texttt{Rakibinin takım kelimeleri: duvar, petunya, beyaz, darp, mağara, hardal, çatı, külot.} \\
\texttt{Dikkat dağıtıcı kelimeler: tavan, tüberküloz, leylak, iglo, oda, kot pantolon, veba.} \\
\texttt{Öldüren kelimeler: hız yapma.} \\
            }
        }
    }
    & \\ \\

    \theutterance \stepcounter{utterance}  
    & \multicolumn{4}{p{0.6\linewidth}}{
        \cellcolor[rgb]{0.9,0.9,0.9}{
            \makecell[{{p{\linewidth}}}]{
                \texttt{\tiny{[P1$\rangle$GM]}}
                \texttt{İPUCU: çiçek} \\
\texttt{HEDEFLER: iris, nergis} \\
            }
        }
    }
    & & \\ \\

    \theutterance \stepcounter{utterance}  
    & & & \multicolumn{2}{p{0.3\linewidth}}{
        \cellcolor[rgb]{0.9,0.9,0.9}{
            \makecell[{{p{\linewidth}}}]{
                \texttt{\tiny{[GM$|$GM]}}
                \texttt{çiçek} \\
            }
        }
    }
    & & \\ \\

    \theutterance \stepcounter{utterance}  
    & & & \multicolumn{2}{p{0.3\linewidth}}{
        \cellcolor[rgb]{0.9,0.9,0.9}{
            \makecell[{{p{\linewidth}}}]{
                \texttt{\tiny{[GM$|$GM]}}
                ['iris', 'nergis']
            }
        }
    }
    & & \\ \\

    \theutterance \stepcounter{utterance}  
    & & & \multicolumn{2}{p{0.3\linewidth}}{
        \cellcolor[rgb]{0.9,0.9,0.9}{
            \makecell[{{p{\linewidth}}}]{
                \texttt{\tiny{[GM$|$GM]}}
                {'word': 'iris', 'assignment': 'team'}
            }
        }
    }
    & & \\ \\

    \theutterance \stepcounter{utterance}  
    & & & \multicolumn{2}{p{0.3\linewidth}}{
        \cellcolor[rgb]{0.9,0.9,0.9}{
            \makecell[{{p{\linewidth}}}]{
                \texttt{\tiny{[GM$|$GM]}}
                {'word': 'nergis', 'assignment': 'team'}
            }
        }
    }
    & & \\ \\

    \theutterance \stepcounter{utterance}  
    & & \multicolumn{4}{p{0.6\linewidth}}{
        \cellcolor[rgb]{0.9,0.9,0.9}{
            \makecell[{{p{\linewidth}}}]{
                \texttt{\tiny{[GM$\rangle$P2]}}
                \texttt{Aşağıdaki listeden, verilen sözcük "çiçek" ile en iyi ilgili veya en yakından bağlantılı, virgülle ayrılmış en fazla 2 sözcükten oluşan bir liste verin. Tahmininizi/tahmini(n)i her zaman TAHMİN: başlayın ve yanıtınıza başka hiçbir metin eklemeyin.} \\
\\ 
\texttt{iglo, hardal, külot, çatı, duvar, şeker, hız yapma, oda, yurt, kot pantolon, tipi, tavan, beyaz, iris, adam kaçırma, darp, tüberküloz, nergis, leylak, gümüş, petunya, veba, mağara, t{-}shirt, etek} \\
            }
        }
    }
    & \\ \\

    \theutterance \stepcounter{utterance}  
    & & & \multicolumn{4}{p{0.6\linewidth}}{
        \cellcolor[rgb]{0.9,0.9,0.9}{
            \makecell[{{p{\linewidth}}}]{
                \texttt{\tiny{[GM$\langle$P2]}}
                \texttt{TAHMİN: nergis, leylak} \\
            }
        }
    }
     \\ \\

    \theutterance \stepcounter{utterance}  
    & & & \multicolumn{2}{p{0.3\linewidth}}{
        \cellcolor[rgb]{0.9,0.9,0.9}{
            \makecell[{{p{\linewidth}}}]{
                \texttt{\tiny{[GM$|$GM]}}
                ['nergis', 'leylak']
            }
        }
    }
    & & \\ \\

    \theutterance \stepcounter{utterance}  
    & & & \multicolumn{2}{p{0.3\linewidth}}{
        \cellcolor[rgb]{0.9,0.9,0.9}{
            \makecell[{{p{\linewidth}}}]{
                \texttt{\tiny{[GM$|$GM]}}
                {'word': 'nergis', 'assignment': 'team'}
            }
        }
    }
    & & \\ \\

    \theutterance \stepcounter{utterance}  
    & & & \multicolumn{2}{p{0.3\linewidth}}{
        \cellcolor[rgb]{0.9,0.9,0.9}{
            \makecell[{{p{\linewidth}}}]{
                \texttt{\tiny{[GM$|$GM]}}
                {'word': 'nergis', 'assignment': 'team'}
            }
        }
    }
    & & \\ \\

    \theutterance \stepcounter{utterance}  
    & & & \multicolumn{2}{p{0.3\linewidth}}{
        \cellcolor[rgb]{0.9,0.9,0.9}{
            \makecell[{{p{\linewidth}}}]{
                \texttt{\tiny{[GM$|$GM]}}
                {'word': 'leylak', 'assignment': 'innocent'}
            }
        }
    }
    & & \\ \\

    \theutterance \stepcounter{utterance}  
    & & & \multicolumn{2}{p{0.3\linewidth}}{
        \cellcolor[rgb]{0.9,0.9,0.9}{
            \makecell[{{p{\linewidth}}}]{
                \texttt{\tiny{[GM$|$GM]}}
                \texttt{leylak} \\
            }
        }
    }
    & & \\ \\

    \theutterance \stepcounter{utterance}  
    & & & \multicolumn{2}{p{0.3\linewidth}}{
        \cellcolor[rgb]{0.9,0.9,0.9}{
            \makecell[{{p{\linewidth}}}]{
                \texttt{\tiny{[GM$|$GM]}}
                {'word': 'çatı', 'assignment': 'opponent'}
            }
        }
    }
    & & \\ \\

    \theutterance \stepcounter{utterance}  
    & & & \multicolumn{2}{p{0.3\linewidth}}{
        \cellcolor[rgb]{0.9,0.9,0.9}{
            \makecell[{{p{\linewidth}}}]{
                \texttt{\tiny{[GM$|$GM]}}
                {'hidden': {'team': ['gümüş', 'iris', 'etek', 'yurt', 'şeker', 'adam kaçırma', 'tipi', 't-shirt'], 'innocent': ['tavan', 'tüberküloz', 'iglo', 'oda', 'kot pantolon', 'veba'], 'opponent': ['duvar', 'petunya', 'beyaz', 'darp', 'mağara', 'hardal', 'külot'], 'assassin': ['hız yapma']}, 'revealed': {'team': {'team': ['nergis'], 'innocent': ['leylak'], 'opponent': [], 'assassin': []}, 'opponent': {'team': [], 'innocent': [], 'opponent': ['çatı'], 'assassin': []}}}
            }
        }
    }
    & & \\ \\

    \theutterance \stepcounter{utterance}  
    & & \multicolumn{4}{p{0.6\linewidth}}{
        \cellcolor[rgb]{0.9,0.9,0.9}{
            \makecell[{{p{\linewidth}}}]{
                \texttt{\tiny{[P1$\langle$GM]}}
                \texttt{nergis kelimesi doğru tahmin edildi. leylak kelimesi tahmin edildi, ancak bu bir masum kelimedir. Takım arkadaşının sırası orada bitti.} \\
\texttt{Şimdi, kalan takım kelimelerinden bazılarıyla ilgili başka bir ipucu ve ilgili hedef kelimelerin bir listesini ver. İpucuna 'İPUCU: ' ile başlamayı, yeni bir satıra geçmeyi ve virgülle ayrılmış hedef kelime listene 'HEDEFLER: ' ile başlamayı unutma. Dikkat: önceki isteklere kıyasla listelerden bazı kelimeler çıkarılmış olabilir.} \\
\\ 
\texttt{Kalan takım kelimelerin: gümüş, iris, etek, yurt, şeker, adam kaçırma, tipi, t{-}shirt.} \\
\texttt{Rakibin için kalan kelimeler: duvar, petunya, beyaz, darp, mağara, hardal, çatı, külot.} \\
\texttt{Kalan dikkat dağıtıcı kelimeler: tavan, tüberküloz, iglo, oda, kot pantolon, veba.} \\
\texttt{Kalan öldüren kelimeler: hız yapma.} \\
            }
        }
    }
    & \\ \\

    \theutterance \stepcounter{utterance}  
    & \multicolumn{4}{p{0.6\linewidth}}{
        \cellcolor[rgb]{0.9,0.9,0.9}{
            \makecell[{{p{\linewidth}}}]{
                \texttt{\tiny{[P1$\rangle$GM]}}
                \texttt{İPUCU: giysi} \\
\texttt{HEDEFLER: etek, t{-}shirt} \\
            }
        }
    }
    & & \\ \\

    \theutterance \stepcounter{utterance}  
    & & & \multicolumn{2}{p{0.3\linewidth}}{
        \cellcolor[rgb]{0.9,0.9,0.9}{
            \makecell[{{p{\linewidth}}}]{
                \texttt{\tiny{[GM$|$GM]}}
                \texttt{giysi} \\
            }
        }
    }
    & & \\ \\

    \theutterance \stepcounter{utterance}  
    & & & \multicolumn{2}{p{0.3\linewidth}}{
        \cellcolor[rgb]{0.9,0.9,0.9}{
            \makecell[{{p{\linewidth}}}]{
                \texttt{\tiny{[GM$|$GM]}}
                ['etek', 't-shirt']
            }
        }
    }
    & & \\ \\

    \theutterance \stepcounter{utterance}  
    & & & \multicolumn{2}{p{0.3\linewidth}}{
        \cellcolor[rgb]{0.9,0.9,0.9}{
            \makecell[{{p{\linewidth}}}]{
                \texttt{\tiny{[GM$|$GM]}}
                {'word': 'etek', 'assignment': 'team'}
            }
        }
    }
    & & \\ \\

    \theutterance \stepcounter{utterance}  
    & & & \multicolumn{2}{p{0.3\linewidth}}{
        \cellcolor[rgb]{0.9,0.9,0.9}{
            \makecell[{{p{\linewidth}}}]{
                \texttt{\tiny{[GM$|$GM]}}
                {'word': 't-shirt', 'assignment': 'team'}
            }
        }
    }
    & & \\ \\

    \theutterance \stepcounter{utterance}  
    & & \multicolumn{4}{p{0.6\linewidth}}{
        \cellcolor[rgb]{0.9,0.9,0.9}{
            \makecell[{{p{\linewidth}}}]{
                \texttt{\tiny{[GM$\rangle$P2]}}
                \texttt{nergis kelimesi doğru tahmin edildi. leylak kelimesi tahmin edildi, ancak bu bir masum kelimedir. Senin sıran orada bitti.} \\
\texttt{Şimdi, aşağıdaki kelime listesinden, 'giysi' sözcüğüyle en iyi ilişkili veya en yakından bağlantılı, en az 1 ve en fazla 2 kelimeden oluşan başka bir virgülle ayrılmış liste sağlayın. Cevabınızı TAHMİN: başlamayı unutmayın. Not: bazı kelimeler, önceki isteklerle karşılaştırıldığında listeden kaldırılmış olabilir.} \\
\\ 
\texttt{iglo, hardal, külot, duvar, şeker, hız yapma, oda, yurt, kot pantolon, tipi, tavan, beyaz, iris, adam kaçırma, darp, tüberküloz, gümüş, petunya, veba, mağara, t{-}shirt, etek} \\
            }
        }
    }
    & \\ \\

    \theutterance \stepcounter{utterance}  
    & & & \multicolumn{4}{p{0.6\linewidth}}{
        \cellcolor[rgb]{0.9,0.9,0.9}{
            \makecell[{{p{\linewidth}}}]{
                \texttt{\tiny{[GM$\langle$P2]}}
                \texttt{TAHMİN: külot, etek} \\
            }
        }
    }
     \\ \\

    \theutterance \stepcounter{utterance}  
    & & & \multicolumn{2}{p{0.3\linewidth}}{
        \cellcolor[rgb]{0.9,0.9,0.9}{
            \makecell[{{p{\linewidth}}}]{
                \texttt{\tiny{[GM$|$GM]}}
                ['külot', 'etek']
            }
        }
    }
    & & \\ \\

    \theutterance \stepcounter{utterance}  
    & & & \multicolumn{2}{p{0.3\linewidth}}{
        \cellcolor[rgb]{0.9,0.9,0.9}{
            \makecell[{{p{\linewidth}}}]{
                \texttt{\tiny{[GM$|$GM]}}
                {'word': 'külot', 'assignment': 'opponent'}
            }
        }
    }
    & & \\ \\

    \theutterance \stepcounter{utterance}  
    & & & \multicolumn{2}{p{0.3\linewidth}}{
        \cellcolor[rgb]{0.9,0.9,0.9}{
            \makecell[{{p{\linewidth}}}]{
                \texttt{\tiny{[GM$|$GM]}}
                \texttt{külot} \\
            }
        }
    }
    & & \\ \\

    \theutterance \stepcounter{utterance}  
    & & & \multicolumn{2}{p{0.3\linewidth}}{
        \cellcolor[rgb]{0.9,0.9,0.9}{
            \makecell[{{p{\linewidth}}}]{
                \texttt{\tiny{[GM$|$GM]}}
                {'word': 'duvar', 'assignment': 'opponent'}
            }
        }
    }
    & & \\ \\

    \theutterance \stepcounter{utterance}  
    & & & \multicolumn{2}{p{0.3\linewidth}}{
        \cellcolor[rgb]{0.9,0.9,0.9}{
            \makecell[{{p{\linewidth}}}]{
                \texttt{\tiny{[GM$|$GM]}}
                {'hidden': {'team': ['gümüş', 'iris', 'etek', 'yurt', 'şeker', 'adam kaçırma', 'tipi', 't-shirt'], 'innocent': ['tavan', 'tüberküloz', 'iglo', 'oda', 'kot pantolon', 'veba'], 'opponent': ['petunya', 'beyaz', 'darp', 'mağara', 'hardal'], 'assassin': ['hız yapma']}, 'revealed': {'team': {'team': ['nergis'], 'innocent': ['leylak'], 'opponent': ['külot'], 'assassin': []}, 'opponent': {'team': [], 'innocent': [], 'opponent': ['çatı', 'duvar'], 'assassin': []}}}
            }
        }
    }
    & & \\ \\

    \theutterance \stepcounter{utterance}  
    & & \multicolumn{4}{p{0.6\linewidth}}{
        \cellcolor[rgb]{0.9,0.9,0.9}{
            \makecell[{{p{\linewidth}}}]{
                \texttt{\tiny{[P1$\langle$GM]}}
                \texttt{külot kelimesi tahmin edildi, ancak bu bir rakip kelimedir. Takım arkadaşının sırası orada bitti.} \\
\texttt{Şimdi, kalan takım kelimelerinden bazılarıyla ilgili başka bir ipucu ve ilgili hedef kelimelerin bir listesini ver. İpucuna 'İPUCU: ' ile başlamayı, yeni bir satıra geçmeyi ve virgülle ayrılmış hedef kelime listene 'HEDEFLER: ' ile başlamayı unutma. Dikkat: önceki isteklere kıyasla listelerden bazı kelimeler çıkarılmış olabilir.} \\
\\ 
\texttt{Kalan takım kelimelerin: gümüş, iris, etek, yurt, şeker, adam kaçırma, tipi, t{-}shirt.} \\
\texttt{Rakibin için kalan kelimeler: duvar, petunya, beyaz, darp, mağara, hardal.} \\
\texttt{Kalan dikkat dağıtıcı kelimeler: tavan, tüberküloz, iglo, oda, kot pantolon, veba.} \\
\texttt{Kalan öldüren kelimeler: hız yapma.} \\
            }
        }
    }
    & \\ \\

    \theutterance \stepcounter{utterance}  
    & \multicolumn{4}{p{0.6\linewidth}}{
        \cellcolor[rgb]{0.9,0.9,0.9}{
            \makecell[{{p{\linewidth}}}]{
                \texttt{\tiny{[P1$\rangle$GM]}}
                \texttt{İPUCU: ev} \\
\texttt{HEDEFLER: yurt, duvar, çatı, tavan, oda} \\
            }
        }
    }
    & & \\ \\

    \theutterance \stepcounter{utterance}  
    & & & \multicolumn{2}{p{0.3\linewidth}}{
        \cellcolor[rgb]{0.9,0.9,0.9}{
            \makecell[{{p{\linewidth}}}]{
                \texttt{\tiny{[GM$|$GM]}}
                {'player': 'Cluegiver', 'type': 'target is hallucination', 'utterance': 'İPUCU: ev\nHEDEFLER: yurt, duvar, çatı, tavan, oda', 'message': "Targeted word 'duvar' was not listed, you can only target words provided in the lists.", 'target': 'duvar', 'previous_targets': ['nergis', 'leylak', 'külot'], 'remaining_words': ['iglo', 'hardal', 'şeker', 'hız yapma', 'oda', 'yurt', 'kot pantolon', 'tipi', 'tavan', 'beyaz', 'iris', 'adam kaçırma', 'darp', 'tüberküloz', 'gümüş', 'petunya', 'veba', 'mağara', 't-shirt', 'etek']}
            }
        }
    }
    & & \\ \\

\end{supertabular}
}

\end{document}
